\chapter{Definições}


Este capítulo apresenta os principais conceitos utilizados ao longo deste relatório técnico. A correta compreensão desses termos é essencial para a análise e interpretação dos dados energéticos, bem como para o embasamento das propostas de eficiência. As definições aqui descritas seguem terminologias amplamente adotadas no setor elétrico, especialmente no âmbito do Programa de Eficiência Energética (\ac{PEE}) da \ac{ANEEL}, conforme estabelecido na Resolução Normativa nº 920/2021 \cite{aneel920}.


\begin{description}
    \item[Demanda Contratada:] Potência elétrica máxima kW  acordada entre a unidade consumidora e a distribuidora de energia. Ultrapassagens dessa demanda resultam em penalidades tarifárias.
    
    \item[Consumo Específico:] Relação entre o consumo de energia elétrica (kWh) e uma variável operacional da unidade, como área construída (m²), número de alunos, leitos hospitalares, entre outros. Permite avaliar a eficiência relativa de diferentes unidades.
    
    \item[Fator de Carga:] Índice que expressa a uniformidade da utilização da demanda ao longo do tempo. É calculado pela razão entre a demanda média e a demanda máxima em um período, indicando o aproveitamento da infraestrutura elétrica.
    
    \item[Eficiência Energética:] Capacidade de um sistema ou instalação de reduzir o consumo de energia elétrica sem comprometer os níveis de conforto, produtividade ou desempenho. Envolve o uso racional da energia por meio de tecnologias e práticas mais eficientes.
    
    \item[\ac{RCB}:] Indicador econômico que relaciona o custo total de implantação de uma medida de eficiência energética com os benefícios econômicos gerados, como a economia na fatura de energia.
    
    \item[\ac{FCP}:] Percentual da carga total da unidade que opera simultaneamente no horário de ponta do sistema elétrico. É um parâmetro importante para dimensionamento e análise de impacto tarifário.
    
    \item[Programa de Eficiência Energética (\ac{PEE}):] Programa regulado pela \ac{ANEEL}, que obriga as distribuidoras de energia elétrica a aplicar parte de sua receita operacional líquida em projetos de eficiência energética junto aos consumidores, promovendo o uso racional dos recursos energéticos.
\end{description}
